\documentclass[11pt]{amsart}

\usepackage[T1]{fontenc}
\usepackage{lmodern}
\usepackage{microtype}
\usepackage{mathtools,amssymb,amsthm}
\usepackage[hidelinks]{hyperref}

\emergencystretch=1.5em

\hypersetup{
  pdftitle={A Multilinear Polynomial with Constants Whose Image on M_2(K) Is
            Not a Vector Space}
}

\newtheorem{theorem}{Theorem}
\theoremstyle{remark}
\newtheorem*{remark}{Remark}

\DeclareMathOperator{\rank}{rank}
\newcommand{\Mn}{M_n(K)}
\newcommand{\im}{\widetilde\omega}

\title[An image that is not a vector space]
{A Multilinear Polynomial with Constants Whose Image on $M_2(K)$
 Is Not a Vector Space}

\date{}

\subjclass[2020]{16S50, 15A03, 16R99}
\keywords{L'vov--Kaplansky conjecture, image of a polynomial map, multilinear
polynomial, polynomial with constants, matrix algebra, rank, counterexample}

\begin{document}

\begin{abstract}
The concluding remarks of the e-print \texttt{arXiv:2312.13865v1} of Panja,
Saini and Singh state, as Conjecture~2, that for every field $K$ the image of
every multilinear polynomial $\omega$ in the free product of $M(n,K)$ with the
free algebra on $x_1,\dots,x_m$ is a vector space. Take $n=m=2$ and
$\omega(x_1,x_2)=x_1e_{11}x_2$, a single monomial with one interior
coefficient. Its image is exactly the set of matrices of rank at most one in
$M_2(K)$, so it contains $\omega(I_2,I_2)=e_{11}$ and
$\omega(e_{21},e_{12})=e_{22}$ but not their sum $e_{11}+e_{22}=I_2$; the
printed statement is therefore false. Every entry of the witness is $0$ or
$1$, so it serves over every field. We address only that sentence, in the
version quoted; the published version of the source we were unable to read.
\end{abstract}

\maketitle

\section{The statement}

Let $K$ be a field and let $\Mn\langle x_1,\dots,x_m\rangle$ denote the free
product of the matrix algebra $\Mn$ with the free algebra on
$x_1,\dots,x_m$; its elements are the \emph{polynomials with constants}, and
each such $\omega$ induces a map $\im\colon\Mn^m\to\Mn$ by evaluation.

Section~5 (\emph{Concluding remarks}) of the e-print \cite{PSS}, on page~33,
calls a polynomial \emph{multilinear} ``if it is linear in each variable'',
offers
\[
  x_1A_1x_2A_2x_3+x_2A_4x_1x_3A_5-x_3x_1A_6x_2,
  \qquad A_i\in M(n,K),
\]
as ``a multilinear polynomial in 3 variables in algebra $M(n,K)$'', writes
``Generalizing Theorem B, we present a generalized version of the
L'vov-Kaplansky conjecture'', and states the following.

\begin{quote}
\textbf{Conjecture 2} (\cite{PSS}, p.~33). Let $K$ be a field. For any
multilinear polynomial in $\omega\in M(n,K)\langle x_1,x_2,\dots,x_m\rangle$,
the image $\im(M(n,K))$ is a vector space.
\end{quote}

\noindent
The wording is as printed, including ``polynomial in $\omega$''. By the
displayed example, coefficients standing \emph{between} variables are inside
the quantifier, and no hypothesis of invertibility or nonvanishing on the
$A_i$ appears; the quantifier on $K$ is ``Let $K$ be a field'', although
\cite{PSS} assumes algebraic closure for its theorems. The witness below is
insensitive to the latter point.

That e-print is the only version on arXiv and carries the comment
``Preliminary version; 34 pp;''. A published version appeared as
\emph{Mathematika} \textbf{71} (2025), no.~3, Paper e70031; the DOI returned
HTTP 403 and we were unable to read it, so we make no claim about its wording.
If the published concluding remarks restrict the positions of the
coefficients, then what follows is a correction to the e-print rather than a
counterexample to the published statement. We make no claim about the
L'vov--Kaplansky conjecture, nor about the theorems of \cite{PSS}, nor about
any statement other than the one displayed above.

\section{The counterexample}\label{sec:witness}

Take $n=2$, $m=2$, let $K$ be any field, and let
\[
  \omega(x_1,x_2)=x_1Ax_2,\qquad A=e_{11},
\]
a single monomial with one interior coefficient, of the shape of the example
quoted above. The object, and the two evaluation points that decide the
matter, are these five matrices, written as lists of rows:

\begin{verbatim}
A   = [[1, 0], [0, 0]]
X1  = [[1, 0], [0, 1]]
Y1  = [[1, 0], [0, 1]]
X2  = [[0, 0], [1, 0]]
Y2  = [[0, 1], [0, 0]]
\end{verbatim}

\noindent
All entries are $0$ or $1$, so the same five matrices are defined over every
field through the unique ring map $\mathbb Z\to K$.

\begin{theorem}\label{thm:main}
Let $K$ be a field and $\omega(x_1,x_2)=x_1e_{11}x_2$ in
$M_2(K)\langle x_1,x_2\rangle$. Then $\omega$ is multilinear and
\[
  \im\bigl(M_2(K)\bigr)=\{M\in M_2(K):\rank M\le1\}.
\]
In particular $e_{11}=\omega(X_1,Y_1)$ and $e_{22}=\omega(X_2,Y_2)$ lie in
the image while $e_{11}+e_{22}=I_2$ does not, so the image is not closed
under addition and Conjecture~2 of \cite{PSS} is false.
\end{theorem}

\begin{proof}
Multilinearity is immediate: $\omega$ is a single monomial containing each of
$x_1,x_2$ exactly once, and $(X,Y)\mapsto XAY$ is $K$-bilinear.

Write $e_1=(1,0)^{\mathsf T}$, so that $A=e_1e_1^{\mathsf T}$. Then for all
$X,Y\in M_2(K)$
\begin{equation}\label{eq:fact}
  XAY=(Xe_1)(e_1^{\mathsf T}Y)=u\,v^{\mathsf T},
\end{equation}
where $u$ is the first column of $X$ and $v^{\mathsf T}$ is the first row
of $Y$.
Every such product is an outer product and therefore has rank at most $1$.
Conversely a matrix of rank at most $1$ is $uv^{\mathsf T}$ for some
$u,v\in K^2$ (rank $0$ being $u=0$), and choosing $X$ with first column $u$
and $Y$ with first row $v^{\mathsf T}$ realises it as $XAY$ by
\eqref{eq:fact}. This proves the displayed equality.

For the two witnesses, $X_1=Y_1=I_2$ gives $\omega(I_2,I_2)=A=e_{11}$, and
$X_2=e_{21}$, $Y_2=e_{12}$ give
$\omega(e_{21},e_{12})=e_{21}e_{11}e_{12}=e_{21}e_{12}=e_{22}$. Both have
rank $1$. Their sum is $I_2$, and $\det I_2=1\ne0$ in every field, so
$\rank(e_{11}+e_{22})=2>1$ and $I_2\notin\im(M_2(K))$.
\end{proof}

\begin{remark}
The image contains $0=\omega(0,0)$ and is closed under scalar multiplication,
since $\lambda(XAY)=(\lambda X)AY$; of the vector space axioms only
additivity fails.
\end{remark}

\section{What was checked}

The proof above is by hand and needs no machine. The accompanying program
\texttt{verify.py} uses the Python standard library only, decides nothing in
floating point, and uses no randomness. It parses the five matrices out of the
verbatim display in Section~\ref{sec:witness} rather than carrying its own
copy, and from them re-derives: the two evaluations and all four ranks, over
$\mathbb Q$ and over $\mathbb F_p$ for $p\le11$; the
factorization~\eqref{eq:fact} on every ordered pair of matrix units of $M_n$
for $n\le4$ over $\mathbb Z$ and on all $262{,}144$ ordered pairs of
$M_3(\mathbb F_2)$; multilinearity by exhaustion; and the set equality of
Theorem~\ref{thm:main} over $\mathbb F_2$ and $\mathbb F_3$. The recorded run
\texttt{verify.output.txt} reports \texttt{VERDICT: ALL 89 CHECKS PASS} and
exit status $0$; the checks it reports beyond those just listed concern
statements this note does not make. No enumeration runs over an infinite
field, and none is needed: the proof of Theorem~\ref{thm:main} is uniform in
$K$. The program cannot reach the published version of \cite{PSS}, so the
qualification recorded in Section~1 is not machine checked.

\begin{thebibliography}{9}

\bibitem{PSS}
S.~Panja, P.~Saini, and A.~Singh,
\emph{Images of polynomial maps with constants},
e-print \href{https://arxiv.org/abs/2312.13865v1}{arXiv:2312.13865v1}, 2023,
the version quoted throughout this note. A published version appeared as
Mathematika \textbf{71} (2025), no.~3, Paper e70031,
\href{https://doi.org/10.1112/mtk.70031}{doi:10.1112/mtk.70031}; we were
unable to read it.

\end{thebibliography}

\end{document}
